Text classifiers deployed behind conversational systems receive spelling variants that preserve intent but alter lexical evidence. The operational question is which generated examples deserve a limited final-refit budget. Random and margin-targeted selection use the same class-balanced refit count; Margin-Targeted Typo Augmentation retains candidates with the lowest true-label score gaps.

Let a clean training utterance and intent generate four deterministic candidates through adjacent swap, deletion, insertion, and keyboard-neighbor substitution. A clean classifier scores each candidate using the true-label margin, defined as the score of the correct intent minus the strongest competing intent. For a fixed total budget, equal per-intent quotas prevent common intents from dominating. Within each intent, retain candidates with the lowest true-label margins. Random selection is the matched control; it receives the identical candidate pool, class quotas, and total example count.

The intended Figure 1 is a compact single-column motivation schematic with three vertically aligned regions: one clean utterance branching into four typo glyphs, a central margin-ranking lane, and a final side-by-side comparison of margin-targeted versus random selection under equal class quotas and equal counts. It is not a result plot. Use only short labels and a portrait-safe crop.
